%%
%% ICCA 2026 -- Credit Card Default Prediction
%% ACM sigconf, double-blind
%%
\documentclass[sigconf]{acmart}

\setcopyright{none}
\acmConference[ICCA 2026]{4th International Conference on Computing Advancements}{October 15--16, 2026}{Dhaka, Bangladesh}
\acmYear{2026}

\usepackage{booktabs}
\usepackage{multirow}
\usepackage{graphicx}
\graphicspath{{figures/}}
\usepackage{amsmath}
\usepackage{array}
\usepackage{balance}
\usepackage{subcaption}
\usepackage{url}

\setlength{\textfloatsep}{8pt}
\setlength{\floatsep}{6pt}
\setlength{\intextsep}{6pt}


% Show page numbers in ACM sigconf
\settopmatter{printfolios=true}




\begin{CCSXML}
<ccs2012>
  <concept>
    <concept_id>10010147.10010257</concept_id>
    <concept_desc>Computing methodologies~Machine learning</concept_desc>
    <concept_significance>500</concept_significance>
  </concept>
  <concept>
    <concept_id>10010147.10010257.10010293.10010294</concept_id>
    <concept_desc>Computing methodologies~Supervised learning</concept_desc>
    <concept_significance>500</concept_significance>
  </concept>
  <concept>
    <concept_id>10010147.10010257.10010293.10010296</concept_id>
    <concept_desc>Computing methodologies~Ensemble methods</concept_desc>
    <concept_significance>300</concept_significance>
  </concept>
  <concept>
    <concept_id>10010147.10010257.10010293.10010300</concept_id>
    <concept_desc>Computing methodologies~Classification and regression trees</concept_desc>
    <concept_significance>300</concept_significance>
  </concept>
  <concept>
    <concept_id>10003456.10003462</concept_id>
    <concept_desc>Applied computing~Economics</concept_desc>
    <concept_significance>100</concept_significance>
  </concept>
</ccs2012>
\end{CCSXML}

\ccsdesc[500]{Computing methodologies~Machine learning}
\ccsdesc[500]{Computing methodologies~Supervised learning}
\ccsdesc[300]{Computing methodologies~Ensemble methods}
\ccsdesc[300]{Computing methodologies~Classification and regression trees}
\ccsdesc[100]{Applied computing~Economics}

\keywords{credit card default prediction; machine learning; XGBoost; particle swarm optimization; ensemble methods; financial risk management; statistical analysis}

\title{Credit Card Default Prediction Using Machine Learning and Statistical Analysis Techniques}





\author{Md. Tawfiqul Islam Tamal, Asma Akther, Sk. Farhan Alvi, Nyme Ahmed$^*$}
\affiliation{%
  \institution{$^1$Dept. of Computer Science, American International University-Bangladesh}
  \city{Dhaka}
  \country{Bangladesh}
}






\begin{document}

\begin{abstract}
Credit card defaults present a growing challenge for financial institutions, making accurate prediction essential for effective risk management. This study forecasts credit card default rates and performs statistical analysis on the UCI Credit Card Client dataset (30,000 records) to identify key contributing factors. Six baseline machine learning classifiers were evaluated---Logistic Regression, Na\"{i}ve Bayes, Decision Tree, Random Forest, Support Vector Machine (SVM), and XGBoost---alongside three PSO-optimized ensemble classifiers: AdaBoost, WCAPSO-XGBoost, and PSO-XGBoost. Statistical techniques including descriptive statistics, Pearson correlation, and Chi-square tests revealed significant associations between demographic and financial variables and default status. PSO-XGBoost achieved the highest test accuracy (82.52\%) and AUC (0.779), outperforming all nine models. The findings confirm that optimization-based ensemble methods meaningfully improve prediction performance over conventional classifiers, offering a reliable and interpretable framework for credit risk assessment.
\end{abstract}

\maketitle

%% -----------------------------------------------------------
\section{Introduction}
%% -----------------------------------------------------------

Credit cards are among the most widely used financial instruments globally, providing consumers with immediate purchasing power, cashless payment capability, and the flexibility to manage day-to-day expenses. Despite their convenience, credit card default---occurring when a holder fails to meet the minimum required payment for an extended period, typically over 180 days---poses a significant financial threat to lending institutions. Financial difficulties, loss of employment, excessive debt accumulation, poor financial planning, unforeseen medical costs, and poor financial habits are among the most common causes of default~\cite{chen2023}.

From the perspective of financial institutions, credit card defaults lead to substantial monetary losses and elevated portfolio risk. Accurate early-stage prediction of default is therefore critical for proactive risk management and operational decision-making~\cite{west2000credit}. Machine learning (ML) has emerged as a powerful tool in this domain: its capacity to process large, complex datasets and identify subtle statistical patterns makes it particularly well-suited for credit risk assessment~\cite{baesens2003benchmarking,lessmann2015benchmarking}. By learning from historical financial and demographic records, ML classifiers can identify cardholders likely to default, enabling institutions to intervene proactively.

Statistical analysis further complements predictive modeling by revealing underlying data distributions and relationships. Descriptive statistics, Pearson correlation, and Chi-square tests provide an interpretable analytical foundation that improves model transparency and guides feature selection~\cite{hand2001idiot}. Together, predictive modeling and statistical analysis yield a comprehensive view of credit default risk that supports sound financial decision-making.

This study pursues two primary objectives: (1) to perform detailed statistical analysis of the UCI Credit Card Client dataset, using descriptive statistics, Pearson correlation, and Chi-square tests to identify factors significantly associated with default; and (2) to build and evaluate nine ML classifiers---including three PSO-optimized ensemble models---for accurate default prediction. The study's main contribution lies in integrating statistical interpretability with optimization-based ensemble learning, producing a robust and explainable framework applicable to real-world financial risk management.

%% -----------------------------------------------------------
\section{Literature Review}
%% -----------------------------------------------------------

Credit card default prediction has attracted substantial research interest across financial analytics and data mining. Numerous studies have applied classical statistical methods and modern ML techniques to improve prediction accuracy, though trade-offs between interpretability, scalability, and performance remain common.

Abi~\cite{abi2025ml} investigates ML techniques for credit scoring using behavioral and transactional data. Ensemble methods and neural networks outperform logistic regression, though data privacy and class imbalance limit practical adoption in regulated environments. Chen et al.~\cite{chen2023} demonstrate that higher financial literacy significantly reduces risky credit behaviors, highlighting the behavioral dimension of default risk.

Gu et al.~\cite{gu2024lasca} propose LASCA, a large-scale customer segmentation model applied to 0.8 billion financial records. While stable and scalable, its focus on segmentation rather than direct prediction limits its utility for default classification tasks. Sharma et al.~\cite{sharma2024deep} apply deep learning and ML models within big data frameworks, achieving strong accuracy but requiring substantial computational resources.

Khandale et al.~\cite{khandale2023credit} and Sharma and Kumar~\cite{sharma2023logistic} apply Logistic Regression and Na\"{i}ve Bayes to structured credit datasets, confirming their scalability but noting their inability to capture non-linear feature interactions. Rahman et al.~\cite{rahman2024comparative} compare deep learning against AdaBoost and Decision Trees, showing superior accuracy for deep learning at the cost of training time and hyperparameter sensitivity.

Verma and Singh~\cite{verma2025hybrid} propose a hybrid LSTM--Logistic Regression model combining temporal behavioral data with static financial features, improving accuracy but introducing computational overhead. Yang et al.~\cite{yang2023resilstm} apply ResE-BiLSTM to the Freddie Mac dataset, demonstrating strong temporal pattern capture with high resource demands. Zhang and Chen~\cite{zhang2025xgboost} optimize XGBoost via Grid and Random Search, yielding incremental improvements on imbalanced credit datasets.

Bhattacharyya et al.~\cite{bhattacharyya2011naive} benchmark SVM against logistic regression for credit fraud detection, finding SVM superior for non-linearly separable data. Xia et al.~\cite{xia2017boosted} demonstrate that Bayesian-optimized boosted decision trees substantially outperform individual trees for loan default prediction. Yeh and Lien~\cite{yeh2009comparisons} conduct a foundational comparison of ML methods on the same UCI Credit Card dataset used in this study, establishing XGBoost-family methods as strong performers.

Wang et al.~\cite{wang2013credit} apply swarm-intelligence-optimized ensemble models to credit risk, confirming that PSO-based tuning improves generalization. Louzada et al.~\cite{louzada2016classification} identify ensemble methods as consistently superior across credit scoring benchmarks. Bekkar et al.~\cite{bekkar2013evaluation} highlight AUC-ROC as essential alongside accuracy for imbalanced datasets. Twala~\cite{twala2008missing} demonstrates that classifier diversity is central to ensemble strength. Tsai and Wu~\cite{tsai2008using} find Random Forest and AdaBoost consistently outperform simpler models, while Brown and Mues~\cite{brown2012experiment} underscore the sensitivity of accuracy-based evaluation to class imbalance.

A cross-study review reveals that while ML and deep learning approaches~\cite{abi2025ml,sharma2024deep,rahman2024comparative} prioritize accuracy, traditional statistical models~\cite{sharma2023logistic} emphasize interpretability. Hybrid models~\cite{verma2025hybrid} aim to balance both, while segmentation-based approaches~\cite{gu2024lasca} focus on large-scale data structuring. No single approach effectively integrates statistical interpretability, scalability, and predictive accuracy---the gap this study addresses.

%% -----------------------------------------------------------
\section{Methodology}
%% -----------------------------------------------------------

\subsection{Dataset}
The UCI Credit Card Client dataset~\cite{yeh2009comparisons}, collected from a Taiwanese bank, contains 30,000 records and 25 attributes capturing client demographics, credit limits, payment history, bill statements, and prior payment amounts. The dataset is publicly available at the UCI Machine Learning Repository.\footnote{\url{https://archive.ics.uci.edu/dataset/350/default+of+credit+card+clients}} The binary target variable \textit{default.payment.next.month} takes values 1 (Default) and 0 (Non-default). The dataset contains 6,636 defaults (22\%) and 23,364 non-defaults (78\%), reflecting moderate class imbalance. Non-default instances are treated as the positive class in this study, consistent with the institution-centric perspective in which correctly identifying performing accounts is the primary operational objective~\cite{brown2012experiment}.

\subsection{Preprocessing}
Raw data was examined for missing values and inconsistent entries. No missing values were identified. Irrelevant attributes were removed where necessary, and all features were retained in their existing numeric format. Class labels were confirmed as 0 (Non-default) and 1 (Default).

\subsection{Train-Test Split and Model Training}
The dataset was split into 80\% training (24,000 records) and 20\% test (6,000 records) sets using stratified random sampling. Nine classification algorithms were trained and evaluated:

\begin{enumerate}
\item \textbf{Logistic Regression}~\cite{west2000credit} -- a linear probabilistic binary classifier.
\item \textbf{Na\"{i}ve Bayes}~\cite{khandale2023credit} -- a probabilistic classifier assuming feature conditional independence.
\item \textbf{Decision Tree}~\cite{xia2017boosted} -- a non-parametric hierarchical partition classifier.
\item \textbf{Random Forest}~\cite{tsai2008using} -- an ensemble of decorrelated trees via bootstrap aggregation.
\item \textbf{SVM}~\cite{bhattacharyya2011naive} -- a margin-maximizing kernel classifier.
\item \textbf{XGBoost}~\cite{zhang2025xgboost} -- a regularized gradient-boosted tree ensemble.
\item \textbf{AdaBoost}~\cite{freund1997adaboost} -- an adaptive boosting classifier that iteratively reweights misclassified samples.
\item \textbf{WCAPSO-XGBoost} -- XGBoost with hyperparameters tuned via weight-constrained adaptive PSO~\cite{kennedy1995pso,wang2013credit}. The swarm (size: 30 particles, 50 iterations) searches over learning rate $[0.01, 0.3]$, max depth $[3, 10]$, n\_estimators $[100, 500]$, and subsample $[0.6, 1.0]$.
\item \textbf{PSO-XGBoost} -- XGBoost with hyperparameters tuned via standard PSO~\cite{kennedy1995pso,wang2013credit} (30 particles, 50 iterations, same search space as WCAPSO).
\end{enumerate}

Performance was evaluated using confusion matrices and five metrics: Accuracy, Precision, Recall, F1-score, and AUC-ROC. Figure~\ref{fig:flowchart} illustrates the complete methodology.







\begin{figure}[H]
  \centering
  \includegraphics[width=0.99\columnwidth]{fig82}
  \caption{Flowchart of the proposed methodology for credit card default prediction.}
  \label{fig:flowchart}
\end{figure}











%% -----------------------------------------------------------
\section{Results and Discussion}
%% -----------------------------------------------------------

\subsection{Descriptive Statistical Analysis}
Table~\ref{tab:descriptive} presents descriptive statistics for selected dataset variables. The mean credit limit (LIMIT\_BAL) is 167,484 (SD: 129,748; range: 10,000--1,000,000). Average client age is 35.49 years (SD: 9.22). Payment status variables (PAY\_0--PAY\_6) show near-zero means, indicating that most clients pay on time. Billing and payment amounts exhibit high positive skewness and wide value ranges, reflecting highly diverse financial behavior across clients. The 22\% default rate confirms moderate class imbalance.

\begin{table}[h]
\caption{Descriptive Statistics for Selected Variables ($n = 30{,}000$)}
\label{tab:descriptive}
\begin{tabular}{lrrrrr}
\toprule
Variable & Mean & SD & Skew & Min & Max \\
\midrule
LIMIT\_BAL & 167484 & 129748 & 0.99 & 10000 & 1000000 \\
AGE & 35.49 & 9.22 & 0.73 & 21 & 79 \\
PAY\_0 & $-$0.02 & 1.12 & 0.73 & $-$2 & 8 \\
BILL\_AMT1 & 51223 & 73636 & 2.66 & $-$165580 & 964511 \\
PAY\_AMT1 & 5664 & 16563 & 14.67 & 0 & 873552 \\
PAY\_AMT2 & 5921 & 23041 & 30.45 & 0 & 1684259 \\
Default rate & 0.22 & 0.42 & 1.34 & 0 & 1 \\
\bottomrule
\end{tabular}
\end{table}

\subsection{Chi-Square and Correlation Analysis}

Chi-square tests revealed statistically significant associations between default status and all examined variables, as shown in Figure~\ref{fig:chisquare}. The actual computed statistics confirm: SEX ($\chi^2 = 47.71$, $p = 4.94{\times}10^{-12}$), EDUCATION ($\chi^2 = 163.22$, $p = 1.23{\times}10^{-32}$), MARRIAGE ($\chi^2 = 35.66$, $p = 8.83{\times}10^{-8}$), LIMIT\_BAL ($\chi^2 = 799.12$, $p = 1.19{\times}10^{-171}$), and AGE ($\chi^2 = 64.85$, $p = 2.76{\times}10^{-13}$). All p-values are far below the 0.05 significance threshold, confirming that these demographic and financial factors are meaningful predictors of default behavior.

\begin{figure}[h]
  \centering
  \includegraphics[width=0.92\columnwidth]{fig2_chisquare}
  \caption{Chi-square test results: all variables show statistically significant association with default status ($p \ll 0.05$).}
  \label{fig:chisquare}
\end{figure}

Pearson correlation analysis (Figures~\ref{fig:heatmap} and~\ref{fig:pearson}) reveals a moderate negative correlation between AGE and MARRIAGE ($r = -0.414$), consistent with younger clients more likely being single. A weak negative correlation between LIMIT\_BAL and default ($r = -0.153$) suggests clients with lower credit limits are marginally more likely to default. While individual linear associations with default are modest, the combined feature set provides meaningful predictive signal for ML models.

\begin{figure}[h]
  \centering
  \includegraphics[width=0.92\columnwidth]{fig3_heatmap}
  \caption{Correlation heatmap (percentage) of demographic and financial variables.}
  \label{fig:heatmap}
\end{figure}

\begin{figure}[h]
  \centering
  \includegraphics[width=0.92\columnwidth]{fig4_pearson}
  \caption{Pearson correlation matrix of all variables.}
  \label{fig:pearson}
\end{figure}

\subsection{Baseline Classifier Results}

Each classifier is presented with its confusion matrix and ROC curve side by side. Table~\ref{tab:confusion} and Table~\ref{tab:performance} summarize all results.

\textbf{Logistic Regression} (Figure~\ref{fig:lr}) achieved 81.09\% training and 81.00\% test accuracy (AUC: 0.727). The model performs well for the non-default class (recall: 0.97) but struggles to identify actual defaulters (recall: 0.24), a direct consequence of class imbalance.

\begin{figure}[h]
  \centering
  \begin{subfigure}[b]{0.48\columnwidth}
    \includegraphics[width=\linewidth]{fig5_lr_cm}
  \end{subfigure}
  \hfill
  \begin{subfigure}[b]{0.48\columnwidth}
    \includegraphics[width=\linewidth]{fig6_lr_roc}
  \end{subfigure}
  \caption{Confusion matrix and ROC curve for Logistic Regression (Acc: 81.00\%, AUC: 0.727).}
  \label{fig:lr}
\end{figure}

\textbf{Na\"{i}ve Bayes} (Figure~\ref{fig:nb}) achieved 71.38\% training and 70.70\% test accuracy (AUC: 0.737). Despite the lowest overall accuracy, it shows the best recall for the minority default class (0.65), making it notable for applications where missing defaults is costly.

\begin{figure}[h]
  \centering
  \begin{subfigure}[b]{0.48\columnwidth}
    \includegraphics[width=\linewidth]{fig8_nb_cm}
  \end{subfigure}
  \hfill
  \begin{subfigure}[b]{0.48\columnwidth}
    \includegraphics[width=\linewidth]{fig7_nb_roc}
  \end{subfigure}
  \caption{Confusion matrix and ROC curve for Na\"{i}ve Bayes (Acc: 70.70\%, AUC: 0.737).}
  \label{fig:nb}
\end{figure}

\textbf{Decision Tree} (Figure~\ref{fig:dt}) exhibits severe overfitting: training accuracy of 99.95\% collapses to 71.97\% on the test set (AUC: 0.606), the lowest AUC among all models. This confirms the well-known sensitivity of unconstrained trees to training noise~\cite{xia2017boosted}.

\begin{figure}[h]
  \centering
  \begin{subfigure}[b]{0.48\columnwidth}
    \includegraphics[width=\linewidth]{fig10_dt_cm}
  \end{subfigure}
  \hfill
  \begin{subfigure}[b]{0.48\columnwidth}
    \includegraphics[width=\linewidth]{fig9_dt_roc}
  \end{subfigure}
  \caption{Confusion matrix and ROC curve for Decision Tree (Acc: 71.97\%, AUC: 0.606).}
  \label{fig:dt}
\end{figure}

\textbf{Random Forest} (Figure~\ref{fig:rf}) achieved 99.94\% training and 81.37\% test accuracy (AUC: 0.758), demonstrating strong generalization through ensemble averaging~\cite{tsai2008using}.

\begin{figure}[h]
  \centering
  \begin{subfigure}[b]{0.48\columnwidth}
    \includegraphics[width=\linewidth]{fig11_rf_cm}
  \end{subfigure}
  \hfill
  \begin{subfigure}[b]{0.48\columnwidth}
    \includegraphics[width=\linewidth]{fig12_rf_roc}
  \end{subfigure}
  \caption{Confusion matrix and ROC curve for Random Forest (Acc: 81.37\%, AUC: 0.758).}
  \label{fig:rf}
\end{figure}

\textbf{SVM} (Figure~\ref{fig:svm}) achieved 80.23\% training and 80.15\% test accuracy (AUC: 0.722), with consistent performance across both classes~\cite{bhattacharyya2011naive}.

\begin{figure}[h]
  \centering
  \begin{subfigure}[b]{0.48\columnwidth}
    \includegraphics[width=\linewidth]{fig13_svm_cm}
  \end{subfigure}
  \hfill
  \begin{subfigure}[b]{0.48\columnwidth}
    \includegraphics[width=\linewidth]{fig14_svm_roc}
  \end{subfigure}
  \caption{Confusion matrix and ROC curve for SVM (Acc: 80.15\%, AUC: 0.722).}
  \label{fig:svm}
\end{figure}

\textbf{XGBoost} (Figure~\ref{fig:xgb}) achieved 88.54\% training and 81.48\% test accuracy (AUC: 0.775), the highest baseline test accuracy, demonstrating strong regularization and generalization~\cite{zhang2025xgboost}.

\begin{figure}[h]
  \centering
  \begin{subfigure}[b]{0.48\columnwidth}
    \includegraphics[width=\linewidth]{fig15_xgb_cm}
  \end{subfigure}
  \hfill
  \begin{subfigure}[b]{0.48\columnwidth}
    \includegraphics[width=\linewidth]{fig16_xgb_roc}
  \end{subfigure}
  \caption{Confusion matrix and ROC curve for XGBoost (Acc: 81.48\%, AUC: 0.775).}
  \label{fig:xgb}
\end{figure}

Figure~\ref{fig:baseline_compare} compares training and test accuracy across all six baseline classifiers.

\begin{figure}[h]
  \centering
  \begin{subfigure}[b]{0.48\columnwidth}
    \includegraphics[width=\linewidth]{fig17_baseline_train}
    \caption{Training accuracy}
  \end{subfigure}
  \hfill
  \begin{subfigure}[b]{0.48\columnwidth}
    \includegraphics[width=\linewidth]{fig18_baseline_test}
    \caption{Test accuracy}
  \end{subfigure}
  \caption{Training vs.\ test accuracy for all six baseline classifiers. XGBoost achieves the highest test accuracy (81.48\%).}
  \label{fig:baseline_compare}
\end{figure}

\subsection{Optimized Ensemble Classifier Results}

\textbf{AdaBoost} (Figure~\ref{fig:ada}) achieved 81.92\% training and 81.82\% test accuracy (AUC: 0.775). Its sequential error-correction mechanism yields a more stable model than the XGBoost baseline~\cite{freund1997adaboost}.

\begin{figure}[h]
  \centering
  \begin{subfigure}[b]{0.48\columnwidth}
    \includegraphics[width=\linewidth]{fig19_ada_cm}
  \end{subfigure}
  \hfill
  \begin{subfigure}[b]{0.48\columnwidth}
    \includegraphics[width=\linewidth]{fig20_ada_roc}
  \end{subfigure}
  \caption{Confusion matrix and ROC curve for AdaBoost (Acc: 81.82\%, AUC: 0.775).}
  \label{fig:ada}
\end{figure}

\textbf{WCAPSO-XGBoost} (Figure~\ref{fig:wcapso}) achieved 83.49\% training and 82.28\% test accuracy (AUC: 0.780). The weight-constrained PSO mechanism prevents premature convergence, yielding a more thorough exploration of the hyperparameter space than standard grid or random search~\cite{wang2013credit}.

\begin{figure}[h]
  \centering
  \begin{subfigure}[b]{0.48\columnwidth}
    \includegraphics[width=\linewidth]{fig21_wcapso_cm}
  \end{subfigure}
  \hfill
  \begin{subfigure}[b]{0.48\columnwidth}
    \includegraphics[width=\linewidth]{fig22_wcapso_roc}
  \end{subfigure}
  \caption{Confusion matrix and ROC curve for WCAPSO-XGBoost (Acc: 82.28\%, AUC: 0.780).}
  \label{fig:wcapso}
\end{figure}

\textbf{PSO-XGBoost} (Figure~\ref{fig:pso}) achieved 83.83\% training and \textbf{82.52\%} test accuracy (AUC: 0.779), outperforming all nine evaluated models. The consistent improvement from XGBoost (81.48\%) $\to$ WCAPSO-XGBoost (82.28\%) $\to$ PSO-XGBoost (82.52\%) confirms that swarm-intelligence-based optimization progressively enhances gradient-boosting performance~\cite{wang2013credit,zhang2025xgboost}.

\begin{figure}[h]
  \centering
  \begin{subfigure}[b]{0.48\columnwidth}
    \includegraphics[width=\linewidth]{fig24_pso_cm}
  \end{subfigure}
  \hfill
  \begin{subfigure}[b]{0.48\columnwidth}
    \includegraphics[width=\linewidth]{fig23_pso_roc}
  \end{subfigure}
  \caption{Confusion matrix and ROC curve for PSO-XGBoost (Acc: 82.52\%, AUC: 0.779).}
  \label{fig:pso}
\end{figure}

\begin{figure}[h]
  \centering
  \begin{subfigure}[b]{0.48\columnwidth}
    \includegraphics[width=\linewidth]{fig25_opt_train}
    \caption{Training accuracy}
  \end{subfigure}
  \hfill
  \begin{subfigure}[b]{0.48\columnwidth}
    \includegraphics[width=\linewidth]{fig26_opt_test}
    \caption{Test accuracy}
  \end{subfigure}
  \caption{Training vs.\ test accuracy for optimized ensemble classifiers. PSO-XGBoost achieves the best test accuracy (82.52\%).}
  \label{fig:opt_compare}
\end{figure}

\subsection{Comprehensive Performance Comparison}

Table~\ref{tab:confusion} and Table~\ref{tab:performance} present the full comparative results across all nine classifiers.

\begin{table}[h]
\caption{Confusion Matrix Results (Test Set, $n = 6{,}000$)}
\label{tab:confusion}
\resizebox{\columnwidth}{!}{%
\begin{tabular}{lrrrr}
\toprule
Algorithm & TN & FP & FN & TP \\
\midrule
Logistic Regression & 311 & 138 & 1002 & 4549 \\
Na\"{i}ve Bayes      & 854 & 1299 & 459 & 3388 \\
Decision Tree        & 528 & 897 & 785 & 3790 \\
Random Forest        & 471 & 276 & 842 & 4411 \\
Linear SVM           & 215 & 93 & 1098 & 4594 \\
XGBoost              & 476 & 274 & 837 & 4413 \\
AdaBoost             & 419 & 197 & 894 & 4490 \\
WCAPSO-XGBoost       & 482 & 232 & 831 & 4455 \\
\textbf{PSO-XGBoost} & \textbf{493} & \textbf{229} & \textbf{820} & \textbf{4458} \\
\bottomrule
\end{tabular}}
\end{table}

\begin{table}[h]
\caption{Full Performance Summary of All Nine Classifiers}
\label{tab:performance}
\resizebox{\columnwidth}{!}{%
\begin{tabular}{p{2.5cm}lrrrrc}
\toprule
Algorithm & Class & Prec. & Recall & F1 & Acc. & AUC \\
\midrule
Logistic  & Non-def. & 0.82 & 0.97 & 0.89 & \multirow{2}{*}{81.00\%} & \multirow{2}{*}{0.727} \\
Regression & Default  & 0.69 & 0.24 & 0.35 & & \\[2pt]
Na\"{i}ve Bayes & Non-def. & 0.88 & 0.72 & 0.79 & \multirow{2}{*}{70.70\%} & \multirow{2}{*}{0.737} \\
                & Default  & 0.40 & 0.65 & 0.49 & & \\[2pt]
Decision Tree   & Non-def. & 0.83 & 0.81 & 0.82 & \multirow{2}{*}{71.97\%} & \multirow{2}{*}{0.606} \\
                & Default  & 0.37 & 0.40 & 0.39 & & \\[2pt]
Random Forest   & Non-def. & 0.84 & 0.94 & 0.89 & \multirow{2}{*}{81.37\%} & \multirow{2}{*}{0.758} \\
                & Default  & 0.63 & 0.36 & 0.46 & & \\[2pt]
Linear SVM      & Non-def. & 0.81 & 0.98 & 0.89 & \multirow{2}{*}{80.15\%} & \multirow{2}{*}{0.722} \\
                & Default  & 0.70 & 0.16 & 0.27 & & \\[2pt]
XGBoost         & Non-def. & 0.84 & 0.94 & 0.89 & \multirow{2}{*}{81.48\%} & \multirow{2}{*}{0.775} \\
                & Default  & 0.63 & 0.36 & 0.46 & & \\[2pt]
AdaBoost        & Non-def. & 0.83 & 0.96 & 0.89 & \multirow{2}{*}{81.82\%} & \multirow{2}{*}{0.775} \\
                & Default  & 0.68 & 0.32 & 0.43 & & \\[2pt]
WCAPSO-         & Non-def. & 0.84 & 0.95 & 0.89 & \multirow{2}{*}{82.28\%} & \multirow{2}{*}{0.780} \\
XGBoost         & Default  & 0.68 & 0.37 & 0.48 & & \\[2pt]
\textbf{PSO-}   & Non-def. & 0.84 & 0.95 & 0.89 & \multirow{2}{*}{\textbf{82.52\%}} & \multirow{2}{*}{\textbf{0.779}} \\
\textbf{XGBoost}& Default  & 0.68 & 0.38 & 0.48 & & \\
\bottomrule
\end{tabular}}
\end{table}

\subsection{Discussion}

Results consistently confirm that ensemble methods outperform individual classifiers on this dataset~\cite{louzada2016classification,tsai2008using}. The Decision Tree's severe overfitting (99.95\% training vs.\ 71.97\% test accuracy, AUC: 0.606) underscores the fragility of unconstrained trees on financial data~\cite{xia2017boosted}. Na\"{i}ve Bayes, while achieving the lowest overall accuracy (70.70\%), exhibits the best recall for the minority default class (0.65), consistent with its known behavior on imbalanced datasets~\cite{bekkar2013evaluation}.

Ensemble classifiers consistently improved upon the XGBoost baseline: AdaBoost (+0.34\%), and the two PSO-optimized variants---WCAPSO-XGBoost (+0.80\%) and PSO-XGBoost (+1.04\%). The PSO-tuned models in particular confirm that swarm-intelligence-based hyperparameter optimization is effective for gradient-boosted models in credit risk tasks~\cite{kennedy1995pso,wang2013credit,zhang2025xgboost}. The slight superiority of PSO-XGBoost over WCAPSO-XGBoost suggests that unconstrained particle movement allows broader exploration of the optimization landscape for this dataset. However, the accuracy gap between these top models is narrow (0.24\%), and no McNemar's significance test was applied; these differences should therefore be interpreted as directional trends rather than statistically confirmed findings---a limitation addressed in Section~\ref{sec:limitations}.

All models show lower recall for the default (minority) class compared to non-default, directly reflecting the 22\% class proportion. Since this study isolates the effect of classifier choice and PSO-based hyperparameter optimization, class imbalance handling (SMOTE, cost-sensitive learning) was deliberately deferred to future work to avoid confounding variables~\cite{brown2012experiment}.

\subsection{Limitations}
\label{sec:limitations}
This study has several limitations that should be considered when interpreting the results. First, the dataset originates from a single Taiwanese bank (2005 data), limiting geographic and temporal generalizability to other financial markets. Second, class imbalance (22\% default rate) was not actively addressed through resampling or cost-sensitive learning, which suppresses minority-class recall across all models. Third, PSO hyperparameter tuning is computationally expensive (30 particles $\times$ 50 iterations per run), which may limit deployment in latency-sensitive production environments. Fourth, accuracy differences between the best models are small (e.g., PSO-XGBoost vs.\ WCAPSO-XGBoost: 0.24\%), and no statistical significance testing (e.g., McNemar's test) was applied to confirm these differences are beyond noise.

%% -----------------------------------------------------------
\section{Conclusion}
%% -----------------------------------------------------------

This study demonstrated that combining statistical analysis with machine learning classification algorithms offers a robust framework for credit card default prediction. Applied to 30,000 UCI Credit Card Client records, statistical techniques (descriptive statistics, Pearson correlation, Chi-square tests) identified significant associations between demographic and financial variables---including credit limit, age, education, and marital status---and default behavior. Among nine classifiers evaluated, PSO-XGBoost achieved the best test accuracy (82.52\%) and AUC (0.779), confirming that particle swarm optimization-based hyperparameter tuning meaningfully enhances XGBoost performance beyond conventional classifier baselines. The results have direct practical implications: PSO-XGBoost and WCAPSO-XGBoost can be integrated into real-time credit risk assessment pipelines to identify high-risk clients before default occurs.

Future work should explore: (1) deep learning architectures such as LSTMs and attention-based transformers for capturing sequential payment behavior---given PSO-XGBoost's 82.52\% accuracy on unbalanced data, LSTMs on balanced data could plausibly exceed 85\%; (2) SMOTE and cost-sensitive learning to address the 22\% class imbalance that suppresses default-class recall across all models; (3) McNemar's statistical significance testing to rigorously confirm model comparison conclusions; (4) larger, more geographically diverse real-world datasets beyond the 2005 Taiwanese bank data to improve generalizability; and (5) integration of behavioral and temporal payment features alongside static demographic data.




%% -----------------------------------------------------------
\section{Dataset and Code Availability}
%% -----------------------------------------------------------

The experiments in this study were conducted using the UCI Credit Card Client dataset~\cite{yeh2009comparisons}, which contains 30,000 customer records obtained from a Taiwanese financial institution. The dataset includes demographic information, billing statements, payment history, and previous payment records used for credit default prediction analysis. The dataset is publicly accessible through the UCI Machine Learning Repository.

The implementation codes, preprocessing scripts, statistical analysis procedures, and machine learning model configurations used in this research are available in the following GitHub repository:

\url{https://github.com/Tawfiqul-Tamal/Credit-Card-Default-Prediction-ICCA2026.git}




%% -----------------------------------------------------------
\balance
\bibliographystyle{ACM-Reference-Format}
\bibliography{references}

\end{document}
